\scanpage{638}


\[
\diagram{
\begin{tikzcd}[ampersand replacement=\&,column sep=.6em,row sep=1.3em,cells={font=\scriptsize}]
\&\&Z_{\rho}\arrow[rrr,hook,"2"]
  \arrow[ddddll,dashed,leftrightarrow,bend right=35,
    "{\unsure{\leftrightarrow}}"'{description}]
\&\&\&\mathcal{N}_{\rho}^{*}\\
\&\&\&Z_{\rho,0}^{!}=\mathcal{M}'(-\bar j)
  \arrow[rrr,hook,"2"]
  \arrow[dl,hook,"2"']
\&\&\&\mathcal{M}'(-\bar j)\mathbin{.}\mu_{\tau'}\;{}^{*}
  \arrow[dl,hook,"2"']\\
\&\&Z_{\rho}^{!}=\mathcal{M}'(-\bar j)\times\mu
  \arrow[uu,hook,"3"']
  \arrow[rrr,hook,"2"]
  \arrow[ddll,hook,dashed,"3"']
\&\&\&\mathcal{N}_{\rho}^{?}
  \arrow[uu,hook,"3"']
  \arrow[ddll,hook,dashed,"3"']\\
\&\&\&\mathcal{M}^{0}(-\bar j)
  \arrow[uu,hook,"2"']
  \arrow[rrr,hook,"2"]
  \arrow[ddl,hook,"2"']
\&\&\&\mathcal{M}'''(-\bar j)
  \arrow[uu,hook,"2"']
  \arrow[ddl,hook,"2"']\\
Z_{\rho}\arrow[rrr,hook,dashed]
\&\&\&\unsure{\mathcal{N}_{\rho}|\mathcal{N}_{\rho}^{*}}
  \arrow[uuuurr,dashed,leftrightarrow,bend left=20,
    "{\unsure{\leftrightarrow}}"'{description}]\\
\&\&\mathcal{M}^{0}(-\bar j)\times\mu
  \arrow[uuu,hook,"2"']
  \arrow[rrr,hook,"2"]
  \arrow[dddll,hook,
    "{\deleted{\unsure{\hookrightarrow}}}"{description}]
\&\&\&\mathcal{M}'''(-\bar j)\times\mu
  \arrow[uuu,hook,"2"']
  \arrow[dddll,hook,
    "{\deleted{\unsure{\hookrightarrow}}}"{description}]\\
\&\&\&1\arrow[uuu,hook]
  \arrow[rrr,hook,"2"]
  \arrow[dl,hook,"2"']
\&\&\&\mu_{\tau'}\arrow[uuu,hook]\arrow[dl,hook,"2"']\\
(38)\&\&\mu\arrow[uu,hook]\arrow[rrr,hook,"2"]
  \arrow[ddll,hook,dashed,"3"']
\&\&\&\mathbb{D}_{\tau'}\arrow[uu,hook]
  \arrow[ddll,hook,dashed,"3"']\\
\deleted{\unsure{L_{\rho}}}
  \arrow[rrr,hook,
    "{\deleted{\unsure{\hookrightarrow}}}"{description}]
\&\&\&\deleted{\unsure{\mathbb{D}_{\rho}}}\\
L_{\rho}\arrow[uuuuu,hook,dashed]
  \arrow[rrr,hook,dashed,"2"]
\&\&\&\mathbb{D}_{\rho}\arrow[uuuuu,hook,dashed]
\end{tikzcd}
}
\]


\margin{$*$ -- Caution,
$\mathcal{M}'(-\bar j)\mathbin{.}\mu_{\tau'}$ is \uline{not}
a subgroup, for $\mu_{\tau'}$ does not normalize
$\mathcal{M}'(-\bar j)$.}

\unsure{a)} Subgroups of finite index in $Z_{\rho}$ -- in the diagram
there are four of them


\[
\diagram{
\begin{tikzcd}[ampersand replacement=\&,column sep=3em,row sep=3em]
\&Z_{\rho}^{!}=\mathcal{M}'(-\bar j)\times\mu
\&\mathcal{M}'(-\bar j)=Z_{\rho,0}^{!}=\mathcal{N}_{\rho,0}^{\unsure{*!}}
  \arrow[l,hook,"2"']\\
(39)\&\mathcal{M}^{0}(-\bar j)\times\mu
  \arrow[u,hook,"2"']
\&\mathcal{M}^{0}(-\bar j)=Z_{\rho,0}^{0!}
  \arrow[u,hook,"2"]
  \arrow[l,hook,"2"]
\end{tikzcd}
}
\]


are they invariant in $Z_{\rho}$? This is clear for the subgroups
\deleted{$\unsure{\mathcal{M}'(-\bar j)}$}
\[
\tag{40}
\left\{
\begin{aligned}
\mathcal{M}'(-\bar j)\times\mu
  &=Z_{\rho}\cap\mathcal{M}^{!}=Z_{\rho}^{!},
  \quad\text{and for}\\
\deleted{\unsure{\mathcal{M}^{0}(-\bar j)\times\mu=}}&\\[-.5ex]
\mathcal{M}^{0}(-\bar j)\times\mu
  &=Z_{\rho}\cap\mathcal{M}^{!}\cap
       \operatorname{Ker}(\varepsilon_{2}:\mathcal{M}\to\mu)
\end{aligned}
\right.
\]
-- this shows that these groups are

\margin{637}

\scanpage{639}
\mspage{623}

even invariant in $\mathcal{N}_{\rho}^{*}$, since
\deleted{$\unsure{\mathcal{M}^{!}}$}.
$Z_{\rho},\mathcal{M}^{!},\operatorname{Ker}\varepsilon_{2}$
are stabilized by $\mathcal{N}_{\rho}^{*}$.
Besides, \deleted{since}
$Z_{\rho,0}^{!}=\mathcal{M}'(-\bar j)$ is
\deleted{of index $2$}\inserted{invariant}
\deleted{\unsure{in}}\inserted{in}
$Z_{\rho}^{!}=\mathcal{M}'(-\bar j)\times\mu$,
\unsure{to} show that it is so in $Z_{\rho}$, it suffices to prove
that it is normalized by $\rho$
(since $Z_{\rho}=L_{\rho}\mathbin{.}Z_{\rho}^{!}$ -- this is
$\pm$ trivial), now let us examine in general for
\deleted{$\unsure{x\in}$} $x\in\mathcal{M}$
whether $x$ normalizes \deleted{$\unsure{\mathcal{M}_{0}^{!}}$}
\[
\mathcal{M}_{0}^{!}\overset{\mathrm{dfn}}{=}
\mathcal{M}(0)\mathbin{.}\Pi_{0},
\]
Since it normalizes $\Pi_{0}$, it suffices to see whether
for $u\in\mathcal{M}(0)$, we have
\[
xux^{-1}\in\mathcal{M}_{0}^{!}
\]
which is also equivalent to
\[
xux^{-1}\equiv u\qquad(\Pi_{0})
\]
i.e.
\[
[x,u]\in\Pi_{0}
\]
now we know that for $x\in\mathcal{G}$,
\inserted{$u\in\mathcal{M}^{!}$,} we have
\[
\tag{*}
\underbrace{[x,u]}_{\inserted{x\,u(x)^{-1}}}^{\unsure{c}}
\equiv
(\operatorname{sg}(x))^{\inserted{\unsure{\xi_{2}(u)|\varepsilon_{2}(u)}}}
\deleted{\unsure{\xi_{2}(u)|\varepsilon_{2}(u)}}
\quad\mathrm{mod}\ \Pi_{0}
\]
where

\[
\diagram{
\begin{tikzcd}[ampersand replacement=\&,column sep=2em]
\operatorname{sg}:\&
\mathcal{G}\arrow[r,"\mathrm{can}"]\&
\mathcal{G}_{3}\arrow[r,"\operatorname{sg}"]\&
\mu.
\end{tikzcd}
}
\]


This shows \deleted{\unsure{that we}} that

\uline{Lemma}
\[
\left\{
\begin{aligned}
&\text{If }x\in\mathcal{G}^{+},\ u\in\mathcal{M}_{0}^{!},
  \quad\text{then }
  \deleted{xux^{-1}\in\mathcal{M}_{0}^{!}}
  \ \text{equivalent conditions}\\
&a)\quad xux^{-1}\in\mathcal{M}_{0}^{!}\\
&b)\quad xux^{-1}\equiv u\quad(\Pi_{0})
       \quad\text{i.e.}\quad [x,u]\in\Pi_{0}\\
&c)\quad .\operatorname{sg}(x)^{
       \inserted{\unsure{\xi_{2}(u)|\varepsilon_{2}(u)}}}
       \deleted{\unsure{\xi_{2}(u)|\varepsilon_{2}(u)}}=1
       \quad\text{i.e.}\quad
       \operatorname{sg}(x)=1\quad\text{\uline{or}}\quad
       \varepsilon_{2}(u)=1
\end{aligned}
\right.
\]

\margin{638}

\scanpage{640}

\deleted{For $x=\rho$, this shows that \unsure{the|the}
\deleted{\unsure{condition}}\inserted{\unsure{first group}}
\unsure{is}}

\uline{Corollary 1} Let $x\in\mathcal{G}^{+\wedge}$.
For $x$ to normalize $\mathcal{M}_{0}^{!}$,
it is necessary and sufficient that we have $\operatorname{sg}(x)=+1$.

\uline{Corollary 2}. The subgroup
$\mathcal{M}_{0}^{\prime\prime!}
=\mathcal{M}''(0)\mathbin{.}\Pi_{0}$
is \deleted{invariant in} normalized by
$\mathcal{G}^{+\wedge}$.

This shows in particular that $\rho$ normalizes
$\mathcal{M}_{0}^{!}$, hence also
$\inserted{\mathcal{M}'(-\bar j)=}
Z_{\rho,0}^{!}=Z_{\rho}\cap\mathcal{M}_{0}^{!}$,
which is therefore invariant in $Z_{\rho}$.
The same therefore holds for its intersection
with $\mathcal{M}^{0}(-\bar j)\times\mu$, i.e.
$\mathcal{M}^{0}(-\bar j)$.
(Or again; this group is equal to
$Z_{\rho}\cap\mathcal{M}_{0}^{\prime\prime!}$,
and $\mathcal{M}_{0}^{\prime\prime!}$ is normalized by $\rho\ldots$).
\uline{Thus the four groups} $(39)$
\uline{are normal in $Z_{\rho}$.} )

To see whether they are normal in
$\unsure{\mathcal{N}_{\rho}^{*}}
=\mu_{\tau'}\mathbin{.}Z_{\rho}$,
it suffices to see whether they are normalized by $\tau'$.
This is clear for $Z_{\rho}^{!}=Z_{\rho}\cap\mathcal{M}^{!}$
-- whence the fact that
$\mathcal{N}_{\rho}^{?}=\mu_{\tau'}\mathbin{.}Z_{\rho}^{!}$
is a subgroup of $\mathcal{N}_{\rho}^{*}$.
To
\deleted{$\mathcal{M}^{0}(-\bar j)=
Z_{\rho}\cap\mathcal{M}_{0}^{\prime\prime!}$, it}
\inserted{see \deleted{\unsure{whether the}} which of the}
three other subgroups
\inserted{considered}
last are normalized, it would therefore be necessary to make explicit
the action of $\tau'$ \deleted{on them.}
on $Z_{\rho}^{!}=\mathcal{M}'(-\bar j)\times\mu$.

\margin{639}

\scanpage{641}
\mspage{623}

Let us note that $\tau\in\mathcal{M}(0)\subset\mathcal{M}_{0}^{!}$
normalizes $\mathcal{M}_{0}^{!}$, so
\deleted{\unsure{one|therefore}}\inserted{it} operates
\deleted{$\unsure{Z}$} on
$\mathcal{M}^{!}
=\underbrace{\mathcal{M}_{0}^{!}}_{\simeq\widetilde{\Gamma}_{\mathbb{Q}}\mathbin{.}\Pi_{0}}
\times\mu$
by the product operation of conjugation by $\tau$ in
$\mathcal{M}_{0}^{!}$, and the identity in $\mu$.
On the other hand we have $\tau'=\tau\sigma$, we have seen that
\[
 (40)\qquad
 \sigma(x)\equiv\varepsilon_{2}(x)x\qquad(\Pi_{0})
 \qquad x\in\mathcal{M}
\]
hence
\[
 (41)\qquad
 \tau'(x)\equiv\varepsilon_{2}(x)\tau(x)\qquad(\Pi_{0})
 \qquad x\in\mathcal{M}
\]
so it follows
\[
 (42)\quad
 \left\{
 \begin{aligned}
 &\text{if }x\in\mathcal{M}_{0}^{!},\quad
   \text{then }\tau'(x)\equiv x\quad(\Pi_{0})
   \inserted{\Longleftrightarrow\tau'(x)\in\mathcal{M}_{0}^{!}}
   \Longleftrightarrow\varepsilon_{2}(x)=1\\
 &\text{in particular, }\mathcal{M}_{0}^{\prime\prime!}
   \text{ is normalized by }\tau
 \end{aligned}
 \right.
\]
Thus \deleted{if $x\in$} applying this to the elements of
$Z_{\rho,0}^{!}=\mathcal{M}'(-\bar j)$ in
$Z_{\rho}^{!}=\mathcal{M}'(-\bar j)\times\mu$,
we find
\[
 (43)\quad
 \left\{
 \begin{aligned}
 &\text{Let }x\in Z_{\rho,0}^{!}=\mathcal{M}'(-\bar j),
   \text{ then }\deleted{\unsure{\tau}}\\
 &\tau'(x)\in Z_{\rho,0}^{!}
   \Longleftrightarrow\varepsilon_{2}(x)=1
 \end{aligned}
 \right.
\]

\uline{Corollary}\quad
$Z_{\rho,0}^{!}\inserted{=\mathcal{M}'(-\bar j)}$
is not normalized by $\tau'$, hence is not invariant in
$\unsure{\mathcal{N}_{\rho}^{*}}$,
but \deleted{the}
$\mathcal{M}^{0}(-\bar j)=Z_{\rho}\cap\mathcal{M}_{0}^{\prime\prime!}$
is normalized by $\tau'$, hence also its product with $\mu$.

Thus the only one of the four subgroups $(39)$ of $Z_{\rho}$
which is not normal in $\mathcal{N}_{\rho}$ is
$\mathcal{M}'(-\bar j)=Z_{\rho,0}^{!}$.
The others give

\margin{640}

\scanpage{642}

rise, by $\frac12$ direct product with $\mu_{\tau'}$,
to \unsure{subgroups|subgroups} of $\mathcal{N}_{\rho}$,
invariant in $\mathcal{N}_{\rho}$.
The smallest of all these subgroups is
$\mathcal{M}^{0}(-\bar j)$, which is of index $24$ in
$\unsure{\mathcal{N}_{\rho}^{*}}$.
The quotient group \deleted{\unsure{is a}}
\[
 \deleted{\unsure{\text{group}}}\qquad
 \mathcal{N}_{\rho}^{*}/\mathcal{M}^{0}(-\bar j)
\]
is therefore a group of order $24$, which moreover receives the group
$\mathbb{D}_{\rho}
\inserted{=\mu_{\tau'}\mathbin{.}L_{\rho}}
\subset\mathcal{N}_{\rho}^{*}$,
normalizer of $L_{\rho}$ in
$\unsure{\operatorname{GL}}(2,\mathbb{Z})$,
a group isomorphic to $D_{6}$, and of order $12$.
The kernel of

\[
\diagram{
 \begin{tikzcd}[ampersand replacement=\&]
 (44)\quad\mathbb{D}_{\rho}
 \arrow[r,hook] \&
 \mathcal{N}_{\rho}^{*}/
 \mathcal{M}^{0}(\bar j)
 \end{tikzcd}
}
\]

is obviously contained in $L_{\rho}=\mathbb{D}_{\rho}\cap Z_{\rho}$
and it is reduced to $1$ for
since we even have
\[
 L_{\rho}\cap Z_{\rho,0}^{!}
 \subset L_{\rho}\cap\mathcal{M}_{0}^{!}
 =L_{\rho}\cap\Pi_{0}=\{1\}.
\]
Thus $\mathbb{D}_{\rho}\inserted{\text{(of order }12\text{)}}$
appears as a subgroup of index $2$ of
$\mathcal{N}_{\rho}^{*}/\mathcal{M}^{0}(-\bar j)$,
just as it is a subgroup of index $2$ of
$\mathcal{G}/\Pi_{0}$. It would be tempting

\margin{641}

\scanpage{643}
\mspage{624}

then to establish an isomorphism
\[
 \mathcal{N}_{\rho}^{*}/\mathcal{M}^{0}(-\bar j)
 \overset{?}{\simeq}\mathcal{G}^{\wedge}/\Pi_{0}
\]
compatible with these embeddings of $\mathbb{D}_{\rho}$.
Let us try it \unsure{this way|also}:

\uline{Proposition}\quad Let us consider the composite homomorphism

\[
\diagram{
 \begin{tikzcd}[ampersand replacement=\&,column sep=large,row sep=large]
 (45)\quad\mathcal{M}\arrow[r,"\Theta_{\mathcal{M}}"]\&
 \operatorname{GL}(2,\widehat{\mathbb{Z}})
 \arrow[r]\&
 \operatorname{GL}(2,\mathbb{Z}/4\mathbb{Z})
 \arrow[r,"\substack{\text{epimorphism of}\\\text{degree }4}"]\&
 \begin{gathered}
 \mathcal{G}^{\unsure{*|\wedge}}/\Pi_{0}\\
 \simeq\mathcal{G}^{\wedge}/\Pi_{0}^{\wedge}
 \end{gathered}\\
 \&\&\underbrace{\operatorname{GL}(2,\mathbb{Z})}_{\mathcal{G}}
 /\mathcal{G}[4]\arrow[u,"\sim"']\&
 \end{tikzcd}
}
\]

\deleted{and the induced homomorphism}
which induces on $\mathcal{G}^{\wedge}\subset\mathcal{M}$
the canonical homomorphism
$\mathcal{G}^{\wedge}\longrightarrow
\mathcal{G}^{\wedge}/\Pi_{0}^{\wedge}$.
Then the homomorphism \deleted{composite}\inserted{induced}

\[
\diagram{
 \begin{tikzcd}[ampersand replacement=\&]
 \mathcal{N}_{\rho}^{*}\arrow[r]\&\mathcal{G}/\Pi_{0}
 \end{tikzcd}
}
\]

has as kernel $\mathcal{M}^{0}(-\bar j)$, and induces by passage
to the quotient an isomorphism

\[
\diagram{
 \begin{tikzcd}[ampersand replacement=\&]
 (46)\quad\mathcal{N}_{\rho}^{*}/\mathcal{M}^{0}(-\bar j)
 \arrow[r,"\sim"]\&\mathcal{G}/\Pi_{0}
 \end{tikzcd}
}
\]


Proof. The elements of $\mathcal{M}^{0}(-\bar j)$
are the elements of the form
\[
 \left\{
 v=\alpha^{-1}u\ \middle|\
 \begin{aligned}
 &u\in\mathcal{M}(0),\quad\alpha\in\Pi_{0}\\
 &u(\rho)=\alpha(\rho)
   \qquad(\text{thus }\varepsilon_{3}(u)=1)\\
 &\varepsilon_{2}(u)=1
 \end{aligned}
 \right\}
\]
they are in $1$-$1$ correspondence with the elements of
\[
 \mathcal{M}^{0}(0)
 =\{u\in\mathcal{M}(0)\mid\varepsilon_{2}(u)=\varepsilon_{3}(u)=1\}.
\]
The image \unsure{of|of the}

\margin{642}